\documentclass[a4paper,10pt]{article}
\usepackage[spanish]{babel}
\usepackage{
			ucs
			,amsfonts
			,amssymb
			,acronym
			,verbatim
			,rotating
			,graphicx
			,enumerate
			,fancyhdr
            ,mathrsfs
            ,array
			}

\usepackage{tikz}
\usetikzlibrary{svg.path}

\pagestyle{fancy}
\addtolength{\textheight}{1.0in}
\addtolength{\textwidth}{0.5in}
\addtolength{\voffset}{-1.0in}

\newcommand{\cat}[1][]{%
    \tikz \fill [scale=1ex/30,yscale=-1,#1] svg "M495.959,867.638c-3.318-4-1.121-7.686,7.098-11.904c7.322-3.76,8.223-4.871,11.365-14.049
		c5.646-16.496,9.951-21.084,24.289-25.893c4.883-1.637,7.111-3.217,6.732-4.771c-3.148-12.889-1.992-16.166,4.48-12.701
		c2.174,1.162,3.924,0.754,6.943-1.619c4.547-3.576,5.105-3.137,6.68,5.264c0.725,3.859,0.445,6.785-0.791,8.273
		c-1.186,1.43-1.496,4.199-0.816,7.295c0.895,4.07,0.186,6.26-3.836,11.861c-6.01,8.367-6.746,18.059-1.752,23.053
		c3.975,3.973,4.066,6.363,0.248,6.363c-3.869,0-8.201-4.926-8.201-9.322c0-4.723-4.635-6.33-7.391-2.564
		c-2.844,3.889-2.738,7.115,0.23,7.115c1.312,0,2.387,1.074,2.387,2.387c0,1.92-2.465,2.385-12.635,2.385
		c-6.947,0-13.047-0.668-13.551-1.484c-0.541-0.875-3.803-0.502-7.951,0.912c-7.744,2.639-11.631,5.996-8.896,7.686
		c0.916,0.566,1.668,1.703,1.668,2.527C502.26,870.757,498.104,870.222,495.959,867.638z";~%
}
%\renewcommand{\cat}{\maltese}




\fancyhf{}
\fancyhead[L]{{\bfseries Programaci\'on I Curso 2012. Pr\'actica 3 \hfill Tipos de Datos - Pascal}}
\fancyfoot[LE,RO]{\thepage}

\fancypagestyle{firstpage}{
\fancyhf{}

\fancyhead[L]{\includegraphics[width=1.75cm]{ucalp.jpg}  \\
\small{UCALP. Facultad de Ciencias Exactas e Ingenier\'ia} \\
{\bfseries Programaci\'on I Curso 2012. Pr\'actica 3 \hfill Tipos de Datos - Pascal}}
\fancyfoot[LE,RO]{\thepage}
}

\thispagestyle{firstpage}

\begin{document}

~\\ \\ \\ \\

\begin{enumerate}

\item Eliminar los par\'entesis innecesarios de las siguientes expresiones.
    \begin{itemize}
        \item[a-] \begin{verbatim} a := ( ( 9 - ( 65 /  7 ) + 60 ) * 26 ) + 40 \end{verbatim}
        \item[b-] \begin{verbatim} b := ( - 65 ) /  ( 7 + ( 60 * 26 ) + 40 ) )\end{verbatim}
    \end{itemize}

\item Calcular en valor de las siguientes expresiones, indicando si en TIPO del resultado obtenido se corresponde con el declarado:
    \begin{itemize}
        \item[]	\verb=VARIABLES:=
            \begin{itemize}
                \item[]	\verb=REALES a, b, c, d, e;=
                \item[] \verb=ENTERAS i, k, l, m;=
                \item[] \verb=LOGICAS p, q, r;=
            \end{itemize}
        \item[] \verb=INICIALIZACIONES:=
            \begin{itemize}
                \item[] \begin{verbatim} n := 0;\end{verbatim}
                \item[] \begin{verbatim} b := 6.75;\end{verbatim}
                \item[] \begin{verbatim} c := 12.3;\end{verbatim}
                \item[] \begin{verbatim} l := 42;\end{verbatim}
            \end{itemize}
    \end{itemize}
\begin{tabular}{p{0.4\textwidth} p{0.4\textwidth}}
    \begin{itemize}
        \item[a-] \begin{verbatim} a := b := c;\end{verbatim}
        \item[c-] \begin{verbatim} i := l + 3 - 193;\end{verbatim}
        \item[e-] \begin{verbatim} q := 4 / 2 <> 0;\end{verbatim}
    \end{itemize}
&
    \begin{itemize}
        \item[b-] \begin{verbatim} k := 6 * 11 - l / 5;\end{verbatim}
        \item[d-] \begin{verbatim} p := ( l + 3 ) >= 10;\end{verbatim}
         \item[f-] \begin{verbatim} r := c + m * 5 - 4 * l / k;\end{verbatim}
    \end{itemize}
\\
\end{tabular}    

\item Considerando las siguientes definiciones de constantes y variables, examine cuales de las sentencias de asignaci\'on son inv\'alidas y cu\'al es el error, en caso contrario indique el resultado de la misma.
    \begin{itemize}
        \item[] \verb=CONSTANTES:=
            \begin{itemize}
                \item[] \begin{verbatim} Pagina = 255;\end{verbatim}
                \item[] \begin{verbatim} Z = `z';\end{verbatim}
                \item[] \begin{verbatim} Gamma = 0.0045;\end{verbatim}
            \end{itemize}
        \item[] \verb=VARIABLES:=
            \begin{itemize}
                \item[] \verb=ENTEROS mx7, mx9;=
                \item[] \verb=REALES tar, rango;=
                \item[] \verb=CARACTERES simbolo;=
                \item[] \verb=LOGICAS respuesta;=
            \end{itemize}
    \end{itemize}

\begin{tabular}{p{0.4\textwidth} p{0.4\textwidth}}
    \begin{itemize}
        \item[a-] \begin{verbatim} mx7 := 924;\end{verbatim}
        \item[c-] \begin{verbatim} rango := Gamma;\end{verbatim}
        \item[e-] \begin{verbatim} rango := respuesta;\end{verbatim}
         \item[g-] \begin{verbatim} tar := mx9;\end{verbatim}
         \item[i-] \begin{verbatim} mx9 := Z;\end{verbatim}
    \end{itemize}
&
    \begin{itemize}
       \item[b-] \begin{verbatim} mx9 := Pagina;\end{verbatim}
       \item[d-] \begin{verbatim} Gamma := tar;\end{verbatim}
        \item[f-] \begin{verbatim} mx7 := Gamma;\end{verbatim}
         \item[h-] \begin{verbatim} simbolo := `t`;\end{verbatim}
          \item[j-] \begin{verbatim} respuesta := VERDADERO;\end{verbatim}
    \end{itemize}
\\
\end{tabular}

\item Asumiendo que \verb=x=, \verb=y=, \verb=z= son variables enteras, ?`qu\'e valor tendr\'an luego de la ejecuci\'on de las siguientes secuencia de instrucciones?\\
\begin{tabular}{p{0.4\textwidth} p{0.4\textwidth}}
    \begin{itemize}
        \item[a-] 
\begin{verbatim} 
x := 50;
y := 340;
z := x + y - 190;
x := 17;
y := x + z;
z := x + 200;
z := x - z - 200;
\end{verbatim}
    \end{itemize}
&
    \begin{itemize}
        \item[b-] 
\begin{verbatim}
x := 27
y := 343;
z := x + y - 300;
x := z / 10 + 23;
y := ( x + z ) / 10 * 10;
x := x + 70 + y;
y := z + x - 70;
\end{verbatim}
    \end{itemize}
\\
\end{tabular}

\item Identificar los errores de las siguientes sentencias de asignaci\'on. Cuando encuentre una sentencia ilegal, indique el error. Si la sentencia es correcta, calcule el valor que se asignar\'a.

\begin{tabular}{p{0.4\textwidth} p{0.4\textwidth}}
    \begin{itemize}
        \item[a-] \begin{verbatim} x1 := b * b / -7 / 1.45;\end{verbatim}
        \item[c-] \begin{verbatim} -v := x * y - total;\end{verbatim}
    \end{itemize}
&
    \begin{itemize}
        \item[b-] \begin{verbatim} x := sqrt( abs( 5 * 5 - 4 * 4));\end{verbatim}
        \item[d-] \begin{verbatim} precioTotal - descuento := precioPagar;\end{verbatim}
    \end{itemize}
\\
\end{tabular}	

\item Escribir la declaraci\'on de las variables correspondientes, al conjunto de sentencias de asignaci\'on que se dan a continuaci\'on; y, calcular el valor si es posible.

\begin{tabular}{p{0.4\textwidth} p{0.4\textwidth}}
    \begin{itemize}
        \item[a-] \begin{verbatim} a := trunc( 3.7 );\end{verbatim}
        \item[c-] \begin{verbatim} c := -round( 78.4 );\end{verbatim}
        \item[e-] \begin{verbatim} e := sqrt( 8 );\end{verbatim}
    \end{itemize}
&
    \begin{itemize}
        \item[b-] \begin{verbatim} b := round( 99.9 );\end{verbatim}
        \item[d-] \begin{verbatim} d := 10 / 4;\end{verbatim}
        \item[f-] \begin{verbatim} f := round( -3.1 ) > trunc( - 3.1 );\end{verbatim}
    \end{itemize}
\\
\end{tabular}

\item[] \begin{large}\textbf{Antes de codificar los siguientes ejercicios, especifique las entradas y salidas requeridas para la resolver el problema.} \end{large}

\item Realizar los programas que devuelvan las siguientes salidas. Utilizando estructuras de control.
\begin{small} \emph{Nota: en la siguiente salida el car\'acter} \verb=_= \emph{es un espacio en blanco.}\end{small}

\begin{tabular}{p{0.2\textwidth} p{0.2\textwidth} p{0.4\textwidth}}
    \begin{itemize}
        \item[a-] 
\begin{verbatim} 
*****
*****
*****
*****
*****
\end{verbatim}
    \end{itemize}
&
    \begin{itemize}
        \item[b-] 
\begin{verbatim} 
_*
__*
___*
____*
_____*
\end{verbatim}
    \end{itemize}
&
    \begin{itemize}
        \item[c-] 
\begin{verbatim}
____*____    
___***___   
__*****__  
_*******_ 
____*____    
\end{verbatim}
    \end{itemize}
\\
\end{tabular}

\item Desarrollar un algoritmo para, dado un conjunto de n\'umeros enteros, escribir aquellos que sean divisibles por $I$.

\item Desarrolle un programa que se le ingrese una cierta cantidad de dinero (en pesos y entera) y devuelva la menor cantidad de billetes de \$100, \$50, \$20, \$10, \$5, \$2 y monedas de \$1 posibles.
    
\item Desarrollar un programa que convierta una cantidad entera de segundos al formato \emph{HORAS:MINUTOS:SEGUNDOS}.

\item Desarrollar un programa que pida ingresar un n\'umero entero por teclado estrictamente mayor a 1 y devuelva su mayor divisor que sea distinto de el.
    
\item Desarrollar un programa que pida ingresar un n\'umero entero por teclado estrictamente mayor a 1 y devuelva todos los divisores de ese n\'umero.
    
\item Desarrolle un programa que permita saber si un n\'umero ingresado por teclado es par (positivo y negativo) o impar (positivo y negativo).

\item Indicar qu\'e escribir\'a cada uno de los siguientes algoritmos. Suponer que todas las variables son del tipo entero.

\begin{tabular}{p{0.4\textwidth} p{0.4\textwidth} }
    \begin{itemize}
        \item[a-]
\begin{verbatim}
var := 0;
for i := 1 to 10 do
begin
    var := var + 1;
    writeln(var);
end;
\end{verbatim}
    \end{itemize}
&
    \begin{itemize}
        \item[b-]
\begin{verbatim}
for k := 1 to 15 step 3 do
        for i := -10 to -5 step 1 do
            var := var + 1;
        end;
end;
\end{verbatim}
    \end{itemize}
\\
\end{tabular}

\item Desarrollar un algoritmo que diga si un n\'umero ingresado es abundante.\\
	\begin{small} \emph{Nota: usar Google o Wikipedia para saber que es un n\'umero abundante.} \end{small}

\item \'Idem anterior, pero para un conjunto de n\'umeros de entrada. El algoritmo para cuando el n\'umero ingresado es igual a 0.

\item Desarrollar un algoritmo para, dado un n\'umero natural, determinar si es defectivo.
	\begin{small} \emph{Nota: usar Google o Wikipedia para saber que es un n\'umero defectivo.} \end{small}

\item Desarrollar un algoritmo para encontrar los n\'umeros perfectos entre 1 y \verb=N=.

\item Codificar los ejercicios \emph{3, 4, 7, 9, 10} de la pr\'actica 2 en PASCAL.

\end{enumerate}
\end{document}
